% Generated draft: requires mathematical and source review.
\documentclass[11pt]{article}
\usepackage[margin=25mm]{geometry}
\usepackage{fontspec}
\setmainfont{Latin Modern Roman}
\usepackage{amsmath,amssymb,mathtools,unicode-math}
\setmathfont{Latin Modern Math}
\usepackage{enumitem,needspace,xurl,hyperref}
\setlength{\parindent}{0pt}
\setlength{\parskip}{5pt}
\setlength{\emergencystretch}{3em}
\newcommand{\N}{\mathbb{N}}
\newcommand{\Z}{\mathbb{Z}}
\newcommand{\Q}{\mathbb{Q}}
\newcommand{\R}{\mathbb{R}}
\newcommand{\C}{\mathbb{C}}
\newcommand{\F}{\mathbb{F}}
\newcommand{\A}{\mathbb{A}}
\newcommand{\PP}{\mathbb{P}}
\newcommand{\E}{\mathbb{E}}
\newcommand{\eps}{\varepsilon}
\newcommand{\dd}{\,\mathrm{d}}
\begin{document}
\section*{UnsolvedMath ID: 700608}
\textbf{Source problem code:} catalogue entry
\subsection{Definitions and mathematical statement}
A formation is a nonempty class \(\mathfrak F\) of finite groups, closed under isomorphisms, homomorphic images, and subdirect products. Here closure under subdirect products means that if \(N_1,N_2\trianglelefteq G\), \(N_1\cap N_2=1\), and \(G/N_i\in\mathfrak F\) for \(i=1,2\), then \(G\in\mathfrak F\).

For a finite group \(X\), let \(\pi(X)\) denote the primes dividing \(|X|\). A chief factor \(A/B\) of \(G\) has \(B<A\) normal in \(G\) with no normal subgroup of \(G\) strictly between them. Define
\[
C_G(A/B)=\{g\in G:[a,g]\in B\text{ for every }a\in A\}.
\]
Under the usual local-definition convention, a formation \(\mathfrak F\) is local if there is a function \(f\) assigning to each prime \(p\) a formation, or the empty class, such that
\[
G\in\mathfrak F
\quad\Longleftrightarrow\quad
G/C_G(A/B)\in f(p)
\quad\text{for every chief factor }A/B\text{ of }G
\text{ and every }p\in\pi(A/B).
\]

For \(K\leq H\), put \(\operatorname{Core}_H(K)=\bigcap_{h\in H}h^{-1}Kh\). Use the standard convention that \(H\leq G\) is \(\mathfrak F\)-subnormal if \(H=G\), or there is a chain
\[
H=H_0<H_1<\cdots<H_n=G
\]
such that \(H_{i-1}\) is maximal in \(H_i\) and
\[
H_i/\operatorname{Core}_{H_i}(H_{i-1})\in\mathfrak F
\qquad(1\leq i\leq n).
\]
Write \(\mathcal S_{\mathfrak F}(G)\) for the set of these subgroups, ordered by inclusion.

\textbf{Classification problem (L. A. Shemetkov, 1984).} Determine all local formations \(\mathfrak F\) of finite groups such that, for every finite group \(G\), \(\mathcal S_{\mathfrak F}(G)\) forms a lattice. Under the literal order-theoretic reading, this means that for every \(H,K\in\mathcal S_{\mathfrak F}(G)\) there exist \(M,J\in\mathcal S_{\mathfrak F}(G)\) satisfying
\[
\begin{gathered}
M\leq H,\quad M\leq K,\qquad
L\leq H,\ L\leq K\ \Longrightarrow\ L\leq M
\quad\text{for every }L\in\mathcal S_{\mathfrak F}(G),\\
H\leq J,\quad K\leq J,\qquad
H\leq U,\ K\leq U\ \Longrightarrow\ J\leq U
\quad\text{for every }U\in\mathcal S_{\mathfrak F}(G).
\end{gathered}
\]
The supplied wording does not specify whether it instead requires a sublattice of the full subgroup lattice, namely
\[
H,K\in\mathcal S_{\mathfrak F}(G)
\quad\Longrightarrow\quad
H\cap K,\ \langle H,K\rangle\in\mathcal S_{\mathfrak F}(G).
\]
These readings must not be identified without justification. The terminology conventions above and the intended lattice reading require source review.

\textbf{Supplied qualification.} A comment dated 2015 states that subgroup-closed formations with the required property are described by Xiaolan Yi and S. F. Kamornikov in the reference below. Subgroup-closed means that \(X\in\mathfrak F\) and \(Y\leq X\) imply \(Y\in\mathfrak F\). This comment concerns the subgroup-closed case; it does not supply a classification for the full requested class.
\subsection{Short English statement}
Classify the local formations of finite groups for which the \(\mathfrak F\)-subnormal subgroups of every finite group form a lattice. The supplied record labels the problem open and reports a 2015 result for subgroup-closed formations. Neither assessment is independently verified here. The intended meaning of lattice, and the source's terminology conventions, require review.
\subsection{Sources}
\begin{itemize}
\item Local UnsolvedMath registry, supplied complete record: permanent ID \(700608\), exact source code \texttt{catalogue entry}, title \emph{Kourovka Notebook Problem 9.75}. The supplied status is \texttt{open}; this is a registry label, not an independently verified mathematical assessment.
\item \emph{The Kourovka Notebook (Unsolved Problems in Group Theory)}, original problem 9.75, ninth issue (1984); proposed by L. A. Shemetkov. The supplied record identifies the twenty-first issue edition (2026), \texttt{arXiv:1401.0300v48}, and names E. I. Khukhro and V. D. Mazurov as editors of the collection.
\item Xiaolan Yi and S. F. Kamornikov, \emph{J. Algebra} \textbf{444} (2015), 143--151. The supplied \emph{Comment of 2015} attributes a description of subgroup-closed formations with the stated property to this reference. No article title or explicit link is supplied; the reference and the scope of the reported result have not been independently checked.
\end{itemize}
\end{document}
